\chapter{C++}\label{ch:iv3:cpp}

\begin{ivtrains}
Lifetime, ownership and undefined behaviour first, then writing the standard library's own classes on a whiteboard: a smart pointer, a vector, an allocator. Stems say what they compile with where the answer depends on it, which here is C++20 on g++ 11.4, libstdc++, x86-64 Linux. Learn what the language promises separately from what this compiler happens to do.
\end{ivtrains}

\begin{question}[$\star$ \iqroles{QD, Dev, LLE} \iqfirm{any} \iqlevel{junior}]\label{iq:iv3:cpp:1}
In \cref{lst:iv3:cp:ptrconst}, replace the comment with each of these lines in turn: \texttt{*p = 3;}, \texttt{p = \&b;}, \texttt{*q = 3;}, \texttt{q = \&b;}, \texttt{*r = 3;}, \texttt{r = \&b;}. Which ones compile? How do you read such a declaration?
\omcode{code/interviews-3/03-cpp/cpp/ptr_const.cpp}{1}{8}{Three const pointers.}\label{lst:iv3:cp:ptrconst}
\end{question}

\begin{question}[$\star$ \iqroles{QD, Dev, LLE} \iqfirm{any} \iqlevel{junior}]\label{iq:iv3:cpp:2}
What does \cref{lst:iv3:cp:storage} print, in order? Which object is never destroyed? Name the storage duration of each of the four objects, and say where each one lives on x86-64 Linux.
\omcode{code/interviews-3/03-cpp/cpp/storage.cpp}{1}{22}{Four kinds of lifetime.}\label{lst:iv3:cp:storage}
\end{question}

\begin{question}[$\star$ \iqroles{QD, Dev, LLE} \iqfirm{any} \iqlevel{junior}]\label{iq:iv3:cpp:3}
In C++, what is the difference between \texttt{class} and \texttt{struct}? In \cref{lst:iv3:cp:classstruct}, replace the comment with each of \texttt{(void)S\{\}.x;}, \texttt{(void)C\{\}.x;}, \texttt{(void)FromC\{\}.get();} and \texttt{(void)FromS\{\}.x;}. Which compile?
\omcode{code/interviews-3/03-cpp/cpp/class_struct.cpp}{1}{8}{Default access.}\label{lst:iv3:cp:classstruct}
\end{question}

\begin{question}[$\star$ \iqroles{QD, Dev, LLE} \iqfirm{any} \iqlevel{junior}]\label{iq:iv3:cpp:4}
What two lines does \cref{lst:iv3:cp:ref} print? What does \texttt{r = b} do, and how does that differ from \texttt{p = \&b}? Give two other differences between a reference and a pointer.
\omcode{code/interviews-3/03-cpp/cpp/ref_rebind.cpp}{1}{13}{A reference and a pointer.}\label{lst:iv3:cp:ref}
\end{question}

\begin{question}[$\star$ \iqroles{QD, Dev, LLE} \iqfirm{any} \iqlevel{junior}]\label{iq:iv3:cpp:5}
With g++ 11.4 on x86-64, with 4-byte \texttt{int} and 8-byte pointers, what three numbers does \cref{lst:iv3:cp:arrayparam} print? Why does the parameter's \texttt{[10]} not help? How do you pass an array and keep its length?
\omcode{code/interviews-3/03-cpp/cpp/array_param.cpp}{2}{12}{An array as a parameter.}\label{lst:iv3:cp:arrayparam}
\end{question}

\begin{question}[$\star$ \iqroles{QD, Dev, LLE} \iqfirm{any} \iqlevel{junior}]\label{iq:iv3:cpp:6}
With g++ on x86-64, what three lines does \cref{lst:iv3:cp:promote} print? What is the type of \texttt{a + b}? Is any of this undefined?
\omcode{code/interviews-3/03-cpp/cpp/promote.cpp}{1}{11}{Small integers.}\label{lst:iv3:cp:promote}
\end{question}

\begin{question}[$\star$ \iqroles{QD, Dev, LLE} \iqfirm{any} \iqlevel{junior}]\label{iq:iv3:cpp:7}
What two lines does \cref{lst:iv3:cp:vectorinit} print? State the rule that decides between the two meanings of each pair of arguments.
\omcode{code/interviews-3/03-cpp/cpp/vector_init.cpp}{1}{8}{Parentheses and braces.}\label{lst:iv3:cp:vectorinit}
\end{question}

\begin{question}[$\star$ \iqroles{QD, Dev, LLE} \iqfirm{any} \iqlevel{junior}]\label{iq:iv3:cpp:8}
What does \cref{lst:iv3:cp:memberorder} print? The initialiser list names \texttt{b} first. What decides the order, and which warning would have told you? Why does it matter if \texttt{b}'s initialiser reads \texttt{a}?
\omcode{code/interviews-3/03-cpp/cpp/member_order.cpp}{2}{19}{The order of members.}\label{lst:iv3:cp:memberorder}
\end{question}

\begin{question}[$\star$ \iqroles{QD, Dev, LLE} \iqfirm{any} \iqlevel{junior}]\label{iq:iv3:cpp:9}
What five lines does \cref{lst:iv3:cp:overload} print? What is the type of \texttt{nullptr}? Why was it added to the language, when 0 already meant a null pointer?
\omcode{code/interviews-3/03-cpp/cpp/overload_zero.cpp}{1}{16}{Zero and nullptr.}\label{lst:iv3:cp:overload}
\end{question}

\begin{question}[$\star$ \iqroles{QD, Dev, LLE, QR} \iqfirm{any} \iqlevel{junior}]\label{iq:iv3:cpp:10}
What does \cref{lst:iv3:cp:removeonly} print? Why does \texttt{std::remove} not change the size of the vector? What is the C++20 one-line way to do both steps?
\omcode{code/interviews-3/03-cpp/cpp/remove_only.cpp}{1}{13}{Removing without erasing.}\label{lst:iv3:cp:removeonly}
\end{question}

\begin{question}[$\star$ \iqroles{QD, Dev, LLE, QR} \iqfirm{any} \iqlevel{junior}]\label{iq:iv3:cpp:11}
What two lines does \cref{lst:iv3:cp:accumulate} print? The code compiles with no warning under \texttt{-Wall -Wextra}. Why is the first number wrong?
\omcode{code/interviews-3/03-cpp/cpp/accumulate_int.cpp}{1}{9}{The starting value of a sum.}\label{lst:iv3:cp:accumulate}
\end{question}

\begin{question}[$\star$ \iqroles{QD, Dev, LLE} \iqfirm{any} \iqlevel{junior}]\label{iq:iv3:cpp:12}
What three lines does \cref{lst:iv3:cp:rangecopy} print? What does \texttt{auto} leave out when it deduces a type from a \texttt{const int\&}? When should a range-based loop use \texttt{const auto\&}, and when is plain \texttt{auto} better?
\omcode{code/interviews-3/03-cpp/cpp/range_copy.cpp}{1}{14}{What auto deduces.}\label{lst:iv3:cp:rangecopy}
\end{question}

\begin{question}[$\star\star$ \iqroles{QD, Dev, LLE} \iqfirm{any} \iqlevel{junior, experienced}]\label{iq:iv3:cpp:13}
\Cref{lst:iv3:cp:bufferbad} compiles with no warning and often seems to run. What is wrong, and what would AddressSanitizer report first? State the rule of five, then write this class to follow it. Could you avoid writing any of those members?
\omcode{code/interviews-3/03-cpp/cpp/buffer_bad.cpp}{1}{14}{A class that owns memory.}\label{lst:iv3:cp:bufferbad}
\end{question}

\begin{question}[$\star\star$ \iqroles{QD, Dev, LLE} \iqfirm{any} \iqlevel{junior, experienced}]\label{iq:iv3:cpp:14}
What two lines does \cref{lst:iv3:cp:nomove} print in C++20? \texttt{Logged} differs from \texttt{Plain} only by an empty destructor, added to help with logging. Which members did that remove? Is the program still correct, and what has it lost? What single line fixes it?
\omcode{code/interviews-3/03-cpp/cpp/no_move.cpp}{1}{26}{An empty destructor.}\label{lst:iv3:cp:nomove}
\end{question}

\begin{question}[$\star\star$ \iqroles{QD, Dev, LLE} \iqfirm{any} \iqlevel{junior, experienced}]\label{iq:iv3:cpp:15}
With libstdc++, whose vector starts empty and doubles its capacity, what two lines does \cref{lst:iv3:cp:vecnoexcept} print? The two types differ only by \texttt{noexcept} on the move constructor. Why does the vector copy in one case? What guarantee is it protecting?
\omcode{code/interviews-3/03-cpp/cpp/vec_noexcept.cpp}{1}{22}{Growth and noexcept.}\label{lst:iv3:cp:vecnoexcept}
\end{question}

\begin{question}[$\star\star$ \iqroles{QD, Dev, LLE} \iqfirm{any} \iqlevel{junior}]\label{iq:iv3:cpp:16}
With libstdc++, what five lines does \cref{lst:iv3:cp:movedfrom} print? Which of those values does the standard promise, and which are only what this library does? What does \texttt{std::move} itself do at run time? What may you do with a moved-from \texttt{std::string}?
\omcode{code/interviews-3/03-cpp/cpp/moved_from.cpp}{1}{21}{What is left after a move.}\label{lst:iv3:cp:movedfrom}
\end{question}

\begin{question}[$\star\star$ \iqroles{QD, Dev, LLE} \iqfirm{any} \iqlevel{junior, experienced}]\label{iq:iv3:cpp:17}
With g++ 11.4 in C++20, what three lines does \cref{lst:iv3:cp:elision} print? Why is writing \texttt{std::move} on the return worse than leaving it out? Why does \texttt{either} copy? And how can \texttt{make} return a type with no usable copy or move constructor?
\omcode{code/interviews-3/03-cpp/cpp/elision.cpp}{2}{35}{Returning a value.}\label{lst:iv3:cp:elision}
\end{question}

\begin{question}[$\star\star$ \iqroles{QD, Dev, LLE} \iqfirm{any} \iqlevel{junior, experienced}]\label{iq:iv3:cpp:18}
What five lines does \cref{lst:iv3:cp:forwarding} print? Why does \texttt{bad} always copy, even when it was handed a temporary? What does \texttt{std::forward<T>} do that \texttt{std::move} does not, and what goes wrong if you use \texttt{std::move} in \texttt{good}?
\omcode{code/interviews-3/03-cpp/cpp/forwarding.cpp}{1}{21}{Handing an argument on.}\label{lst:iv3:cp:forwarding}
\end{question}

\begin{question}[$\star\star$ \iqroles{QD, Dev, LLE} \iqfirm{any} \iqlevel{experienced}]\label{iq:iv3:cpp:19}
What three lines does \cref{lst:iv3:cp:collapse} print? Which of \texttt{T\&\&} in \texttt{f}, \texttt{std::vector<T>\&\&} in \texttt{g}, and \texttt{auto\&\&} are forwarding references, and how can you tell? What is \texttt{T} deduced to in \texttt{f(i)}, and how does \texttt{int\& \&\&} become a legal type? Replace the comment with \texttt{std::vector<int> l; g(l);}: does it compile?
\omcode{code/interviews-3/03-cpp/cpp/collapse.cpp}{1}{23}{Two meanings of a double ampersand.}\label{lst:iv3:cp:collapse}
\end{question}

\begin{question}[$\star\star$ \iqroles{QD, Dev, LLE} \iqfirm{any} \iqlevel{junior, experienced}]\label{iq:iv3:cpp:20}
\Cref{lst:iv3:cp:svtemp} compiles with no warning under g++ 11.4, \texttt{-Wall -Wextra}. What is wrong? What does AddressSanitizer report? Would the bug show with a string of 5 characters? How do you take a \texttt{string\_view} safely, and where is it a good parameter type but a poor member?
\omcode{code/interviews-3/03-cpp/cpp/sv_temp.cpp}{1}{10}{A view of a temporary.}\label{lst:iv3:cp:svtemp}
\end{question}

\begin{question}[$\star\star$ \iqroles{QD, Dev, LLE} \iqfirm{any} \iqlevel{junior}]\label{iq:iv3:cpp:21}
What two lines does \cref{lst:iv3:cp:svcstr} print? A \texttt{string\_view} is passed to a C function through \texttt{.data()}. Why is that a bug, and why does it often pass its tests? Is it undefined here?
\omcode{code/interviews-3/03-cpp/cpp/sv_cstr.cpp}{1}{11}{A view and the end of a string.}\label{lst:iv3:cp:svcstr}
\end{question}

\begin{question}[$\star\star$ \iqroles{QD, Dev, LLE} \iqfirm{any} \iqlevel{junior}]\label{iq:iv3:cpp:22}
What does \cref{lst:iv3:cp:cycle} print? Which line never prints, and why? What does a leak checker say? Fix it with \texttt{std::weak\_ptr}, and say how you use the other node once its pointer is weak.
\omcode{code/interviews-3/03-cpp/cpp/cycle.cpp}{1}{15}{Two nodes that hold each other.\protect\hfil\protect\break}\label{lst:iv3:cp:cycle}
\end{question}

\begin{question}[$\star\star$ \iqroles{QD, Dev, LLE} \iqfirm{any} \iqlevel{junior}]\label{iq:iv3:cpp:23}
Write a simplified \texttt{std::unique\_ptr<T>} in C++20, with no custom deleter: a pointer constructor, destructor, move construction and assignment, no copying, and \texttt{get}, \texttt{release}, \texttt{reset} and the operators. What happens on self-move assignment? Why does \texttt{reset} store the new pointer before deleting the old one? How large is your class?
\end{question}

\begin{question}[$\star\star$ \iqroles{QD, Dev, LLE} \iqfirm{any} \iqlevel{junior, experienced}]\label{iq:iv3:cpp:24}
A type holds a megabyte, and \cref{lst:iv3:cp:weak} replaces \texttt{operator delete} to count frees. With libstdc++, what does it print? When does the megabyte go back in each case, and why do the two differ? Give a case for avoiding \texttt{make\_shared}, and the case that usually favours it.
\omcode{code/interviews-3/03-cpp/cpp/make_shared_weak.cpp}{6}{38}{When the memory goes back.}\label{lst:iv3:cp:weak}
\end{question}

\begin{question}[$\star\star$ \iqroles{QD, Dev, LLE} \iqfirm{any} \iqlevel{junior}]\label{iq:iv3:cpp:25}
What is wrong with \cref{lst:iv3:cp:twoowners}, and what do the two counts show? A method of a class wants to hand out a \texttt{shared\_ptr} to its own object. Why is \texttt{std::shared\_ptr<Order>(this)} the same bug? What is the right way, and what happens if the object is not already owned by a \texttt{shared\_ptr}?
\omcode{code/interviews-3/03-cpp/cpp/two_owners.cpp}{1}{14}{Two owners of one raw pointer.}\label{lst:iv3:cp:twoowners}
\end{question}

\begin{question}[$\star\star$ \iqroles{QD, Dev, LLE} \iqfirm{any} \iqlevel{junior}]\label{iq:iv3:cpp:26}
\Cref{lst:iv3:cp:vecref} may print 9 or 1, or something else. Why? What does AddressSanitizer report, and on which line? Give two fixes, one of which keeps the reference. Does the same code on a \texttt{std::deque} have this bug?
\omcode{code/interviews-3/03-cpp/cpp/vec_ref.cpp}{1}{10}{A reference into a vector.}\label{lst:iv3:cp:vecref}
\end{question}

\begin{question}[$\star\star$ \iqroles{QD, Dev, LLE} \iqfirm{any} \iqlevel{junior, experienced}]\label{iq:iv3:cpp:27}
With libstdc++, what does \cref{lst:iv3:cp:vecbool} print? Why do \texttt{auto} on a \texttt{vector<int>} and on a \texttt{vector<bool>} behave differently? Does \texttt{bool\& b = v[0];} compile? Two threads each write a different index of a \texttt{vector<bool>}: is that a data race, and is it one for \texttt{vector<char>}?
\omcode{code/interviews-3/03-cpp/cpp/vec_bool.cpp}{1}{13}{A vector unlike the others.}\label{lst:iv3:cp:vecbool}
\end{question}

\begin{question}[$\star\star$ \iqroles{QD, Dev, LLE} \iqfirm{any} \iqlevel{junior, experienced}]\label{iq:iv3:cpp:28}
With g++ 11.4 and libstdc++, what three lines does \cref{lst:iv3:cp:stability} print? Which results does the standard promise? What does a rehash invalidate in a \texttt{std::unordered\_map}? You keep pointers to orders held inside a container: which containers allow it?
\omcode{code/interviews-3/03-cpp/cpp/stability.cpp}{1}{24}{Does an element stay put?}\label{lst:iv3:cp:stability}
\end{question}

\begin{question}[$\star\star$ \iqroles{QD, Dev, LLE} \iqfirm{any} \iqlevel{junior, experienced}]\label{iq:iv3:cpp:29}
What three lines does \cref{lst:iv3:cp:emplace} print? When is \texttt{emplace\_back} actually better than \texttt{push\_back}? Replace the comment with \texttt{rows.push\_back(1000000);}: what happens, and why is that better than what \texttt{emplace\_back} did on the line above?
\omcode{code/interviews-3/03-cpp/cpp/emplace.cpp}{1}{22}{Two ways to add an element.}\label{lst:iv3:cp:emplace}
\end{question}

\begin{question}[$\star\star$ \iqroles{QD, Dev, LLE, QR} \iqfirm{any} \iqlevel{junior, experienced}]\label{iq:iv3:cpp:30}
\Cref{lst:iv3:cp:sortle} sorts 100 equal values with a comparator written as \texttt{a <= b}. With g++ 11.4, what happens under AddressSanitizer? What rule did the comparator break? Why might the same code pass on short or distinct inputs? How would you write a comparator on two fields?
\omcode{code/interviews-3/03-cpp/cpp/sort_le.cpp}{1}{9}{A comparator with an equals sign.}\label{lst:iv3:cp:sortle}
\end{question}

\begin{question}[$\star\star$ \iqroles{QD, Dev, LLE} \iqfirm{any} \iqlevel{junior, experienced}]\label{iq:iv3:cpp:31}
With g++ 11.4, \texttt{-std=c++20 -O2}, x86-64, what does \cref{lst:iv3:cp:overflow} print? The first function is a common attempt to detect overflow, for $b > 0$. Why is it wrong, and what does the compiler do with it? What do UBSan and \texttt{-fwrapv} change? Write a portable check that does not use a built-in.
\omcode{code/interviews-3/03-cpp/cpp/overflow_check.cpp}{1}{15}{Checking for overflow too late.}\label{lst:iv3:cp:overflow}
\end{question}

\begin{question}[$\star\star$ \iqroles{QD, Dev, LLE} \iqfirm{any} \iqlevel{junior, experienced}]\label{iq:iv3:cpp:32}
With g++ 11.4, C++20, \texttt{-O2}, x86-64, run with no arguments, what four lines does \cref{lst:iv3:cp:shiftabs} print? For each line, is the value promised by the standard, or undefined? What does the shift print and why? Which of them would UBSan flag?
\omcode{code/interviews-3/03-cpp/cpp/shift_abs.cpp}{1}{13}{Four integer edge cases.}\label{lst:iv3:cp:shiftabs}
\end{question}

\begin{question}[$\star\star$ \iqroles{QD, Dev, LLE} \iqfirm{any} \iqlevel{junior, experienced}]\label{iq:iv3:cpp:33}
In \cref{lst:iv3:cp:lifetime}, which of the references \texttt{one}, \texttt{two} and \texttt{three} are safe to read after their declaration? State the rule. Why does passing the temporary through a function that returns its own argument defeat it? What would AddressSanitizer report for the last one?
\omcode{code/interviews-3/03-cpp/cpp/lifetime.cpp}{1}{15}{Binding a temporary to a reference.}\label{lst:iv3:cp:lifetime}
\end{question}

\begin{question}[$\star\star$ \iqroles{QD, Dev, LLE} \iqfirm{any} \iqlevel{junior, experienced}]\label{iq:iv3:cpp:34}
What four lines does \cref{lst:iv3:cp:defaultargs} print? Why does the first say \texttt{Derived 1}? Why does \texttt{d.g(1)} convert 1 to a \texttt{double} when \texttt{Base} has a better match? Give the fix for each, and say which one a compiler option can warn about.
\omcode{code/interviews-3/03-cpp/cpp/default_args.cpp}{1}{21}{Defaults and hiding.}\label{lst:iv3:cp:defaultargs}
\end{question}

\begin{question}[$\star\star$ \iqroles{QD, Dev, LLE} \iqfirm{any} \iqlevel{junior, experienced}]\label{iq:iv3:cpp:35}
With g++ 11.4, libstdc++ and x86-64, what three lines does \cref{lst:iv3:cp:sizes} print? Why is an empty struct not of size 0, and why is \texttt{Base1} smaller than \texttt{Member}? Why does a lambda deleter cost nothing, and a function pointer deleter cost 8 bytes? What does the 16 of \texttt{shared\_ptr} hold?
\omcode{code/interviews-3/03-cpp/cpp/sizes.cpp}{1}{23}{What things occupy.}\label{lst:iv3:cp:sizes}
\end{question}

\begin{question}[$\star\star$ \iqroles{QD, Dev, LLE} \iqfirm{any} \iqlevel{junior}]\label{iq:iv3:cpp:36}
A header declares \texttt{template <class T> T total(const std::vector<T>\& v);}. Its body is in \texttt{sum.cpp}, and \texttt{main.cpp} includes the header and calls \texttt{total} on a \texttt{std::vector<int>}. Each file compiles, and then the build fails. At which stage, and why? Give the two fixes. What does the second one cost a user who wants \texttt{std::vector<long>}?
\end{question}

\begin{question}[$\star\star$ \iqroles{QD, Dev, LLE} \iqfirm{any} \iqlevel{junior, experienced}]\label{iq:iv3:cpp:37}
What two lines does \cref{lst:iv3:cp:concepts} print? How does the compiler choose between the two versions of \texttt{kind}? Replace the comment with \texttt{notional(1, 2);}: what happens, and where does the error point? How would you have written this constraint before C++20?
\omcode{code/interviews-3/03-cpp/cpp/concepts.cpp}{1}{25}{Constraining a template.}\label{lst:iv3:cp:concepts}
\end{question}

\begin{question}[$\star\star$ \iqroles{QD, Dev, LLE} \iqfirm{crypto, market maker, hedge fund} \iqlevel{junior, experienced}]\label{iq:iv3:cpp:38}
Market data messages are one of three types, and \cref{lst:iv3:cp:variant} holds them in a \texttt{std::variant}. With g++ 11.4 on x86-64, what four lines does it print? How does this compare with a base class and three derived classes, for allocation and for adding a new type? If a new alternative's constructor throws during \texttt{emplace}, can a variant end up holding nothing?
\omcode{code/interviews-3/03-cpp/cpp/msg_variant.cpp}{1}{31}{Messages in a variant.}\label{lst:iv3:cp:variant}
\end{question}

\begin{question}[$\star\star$ \iqroles{QD, Dev, LLE} \iqfirm{any} \iqlevel{junior, experienced}]\label{iq:iv3:cpp:39}
You want a \texttt{std::unordered\_map} keyed by a struct of a symbol and a venue. What two things must you supply, and what happens if you supply neither? \Cref{lst:iv3:cp:keyhash} fills it with 100 symbols on 50 venues each, under two hash functions. With libstdc++, what is the largest bucket under the weak hash, and roughly how large is it under the mixed one? Does the weak hash give wrong answers? What does it cost?
\omcode{code/interviews-3/03-cpp/cpp/key_hash.cpp}{1}{36}{A hash for a struct.}\label{lst:iv3:cp:keyhash}
\end{question}

\begin{question}[$\star\star$ \iqroles{QD, Dev} \iqfirm{any} \iqlevel{junior}]\label{iq:iv3:cpp:40}
Value types, seen from another language. With .NET 8, what four lines does \cref{lst:iv3:cp:boxing} print? \texttt{Fill} is a C\# \texttt{struct}. What is boxing, when does it happen here, and why does the \texttt{Add} through the interface not change \texttt{a}? \Cref{lst:iv3:cp:valuecpp} is the C++ version, which prints \texttt{3 1 9}: what does C++ have instead of boxing?
\omcode{code/interviews-3/03-cpp/csharp/Boxing/Program.cs}{1}{24}{A struct in C\#.}\label{lst:iv3:cp:boxing}
\omcode{code/interviews-3/03-cpp/cpp/value_cpp.cpp}{1}{15}{The same in C++.}\label{lst:iv3:cp:valuecpp}
\end{question}


\begin{question}[$\star\star\star$ \iqroles{QD, Dev, LLE} \iqfirm{any} \iqlevel{junior, experienced}]\label{iq:iv3:cpp:41}
\texttt{Item} has a copy constructor and a \texttt{noexcept} move constructor, and each adds one to \texttt{Item::moves}. With libstdc++, what two lines does \cref{lst:iv3:cp:growth} print? Why does a vector grow by a factor, when growing by 16 elements each time would waste less memory? Give the total number of moves for 1000 insertions under each policy. Why might a library choose a factor below 2? Then \cref{lst:iv3:cp:naivevec}: what does \texttt{v.push\_back(v[0])} do, and how must the real one avoid it?
\omcode{code/interviews-3/03-cpp/cpp/growth.cpp}{11}{22}{Watching the capacity.\protect\hfil\protect\break}\label{lst:iv3:cp:growth}
\omcode{code/interviews-3/03-cpp/cpp/naive_vec.cpp}{6}{37}{A first attempt at a vector.}\label{lst:iv3:cp:naivevec}
\end{question}

\begin{question}[$\star\star\star$ \iqroles{QD, Dev, LLE} \iqfirm{any} \iqlevel{experienced}]\label{iq:iv3:cpp:42}
Write a simplified \texttt{std::optional<T>} in C++20, with \texttt{emplace}, \texttt{reset}, a conversion to \texttt{bool}, \texttt{operator*}, and copy and move. \texttt{T} need not have a default constructor. Why can it not simply have a member \texttt{T value}? What must the storage look like? At which moment do you set the flag, and why? \Cref{lst:iv3:cp:misaligned} places a \texttt{double} at an odd address. x86-64 loads it without complaint, so is the program correct?
\omcode{code/interviews-3/03-cpp/cpp/misaligned.cpp}{4}{8}{An object at the wrong address.}\label{lst:iv3:cp:misaligned}
\end{question}

\begin{question}[$\star\star\star$ \iqroles{QD, Dev, LLE} \iqfirm{crypto, market maker, hedge fund} \iqlevel{experienced}]\label{iq:iv3:cpp:43}
A trading system creates and destroys many orders of one type, and must not touch the heap once running. In C++20, write a fixed-size pool \texttt{Pool<T, N>}: \texttt{create(args.{}.{}.)} returns a \texttt{T*}, or \texttt{nullptr} when full, and \texttt{destroy(p)} takes it back, both in constant time. Where do you keep the list of free slots? What should happen if \texttt{T}'s constructor throws? What mistakes does your pool not catch?
\end{question}

\begin{question}[$\star\star\star$ \iqroles{QD, Dev, LLE} \iqfirm{any} \iqlevel{experienced}]\label{iq:iv3:cpp:44}
In C++20, write a bump arena over a fixed buffer, and an allocator type that lets \texttt{std::vector<int>} and \texttt{std::map<int, double>} take all their memory from it. What must an allocator provide at the minimum? Why does it need a converting constructor from an allocator of another type, when your code only ever names \texttt{pair}? What does \texttt{deallocate} do, and what happens to the old blocks as a vector grows?
\end{question}

\begin{question}[$\star\star\star$ \iqroles{QD, Dev, LLE} \iqfirm{any} \iqlevel{experienced}]\label{iq:iv3:cpp:45}
With g++ 11.4 and libstdc++, what three lines does \cref{lst:iv3:cp:pmr} print? The program replaces the global \texttt{operator new} so that it adds 1 to \texttt{news}. Explain 1020 and 400. Why do the strings in the second block use the buffer, though nobody passed it to them? What does this design gain over an allocator passed as a template argument, and what does it cost?
\omcode{code/interviews-3/03-cpp/cpp/pmr_use.cpp}{19}{57}{Polymorphic memory resources. Here \texttt{pmr} is an alias of \texttt{std::pmr}.\protect\hfil\protect\break}\label{lst:iv3:cp:pmr}
\end{question}

\begin{question}[$\star\star\star$ \iqroles{QD, Dev, LLE} \iqfirm{any} \iqlevel{junior, experienced}]\label{iq:iv3:cpp:46}
In \cref{lst:iv3:cp:assignbad}, \texttt{get} stands for an allocation that can throw. Find two separate faults in \texttt{operator=}. For each, say exactly what happens: \texttt{a = a}, and \texttt{a = b} when \texttt{get} throws. Then rewrite the operator so that both are correct. What does your version cost, and when might you choose not to pay it?
\omcode{code/interviews-3/03-cpp/cpp/assign_bad.cpp}{1}{27}{An assignment with two faults.}\label{lst:iv3:cp:assignbad}
\end{question}

\begin{question}[$\star\star\star$ \iqroles{QD, Dev, LLE} \iqfirm{crypto, market maker, hedge fund} \iqlevel{experienced}]\label{iq:iv3:cpp:47}
With g++ 11.4, \texttt{-std=c++20 -O2}, x86-64, what two lines does \cref{lst:iv3:cp:aliasing} print? Why is the first line not 0, when the store through \texttt{b} really wrote zeros over \texttt{x}? What does \texttt{-fno-strict-aliasing} change? You are parsing a 4-byte price from a packet in a \texttt{const char*} buffer. How do you read it correctly, and what does that cost?
\omcode{code/interviews-3/03-cpp/cpp/aliasing.cpp}{1}{15}{One address, two types.}\label{lst:iv3:cp:aliasing}
\end{question}

\begin{question}[$\star\star\star$ \iqroles{QD, Dev, LLE} \iqfirm{any} \iqlevel{experienced}]\label{iq:iv3:cpp:48}
\Cref{lst:iv3:cp:nullcheck} has two functions that differ only in the order of two lines. With g++ 11.4, \texttt{-O2}, x86-64, what does each one compile to? Why was a null check removed from one of them? What happens on \texttt{read\_bad(nullptr)}? Which option keeps the check, and is that the right fix?
\omcode{code/interviews-3/03-cpp/cpp/nullcheck.cpp}{1}{10}{A check placed after the use.}\label{lst:iv3:cp:nullcheck}
\end{question}

\begin{question}[$\star\star\star$ \iqroles{QD, Dev, LLE} \iqfirm{any} \iqlevel{junior, experienced}]\label{iq:iv3:cpp:49}
\Cref{lst:iv3:cp:maxdangle} compiles with no warning under g++ 11.4, \texttt{-Wall -Wextra}. Run it with no arguments, then with three. Under AddressSanitizer, what happens each time, and why does the number of arguments matter? Why does \texttt{std::max} return a reference at all? Why does \texttt{std::max(1, 2.0)} not compile? Give the fix for this listing.
\omcode{code/interviews-3/03-cpp/cpp/max_dangle.cpp}{1}{7}{The larger of two.}\label{lst:iv3:cp:maxdangle}
\end{question}

\begin{question}[$\star\star\star$ \iqroles{QD, Dev, LLE} \iqfirm{crypto, market maker, hedge fund} \iqlevel{experienced}]\label{iq:iv3:cpp:50}
A live configuration is held in a \texttt{std::shared\_ptr}: one thread replaces it, and another copies it to read it. \texttt{shared\_ptr}'s reference count is atomic. In \cref{lst:iv3:cp:sharedrace}, run with argument \texttt{A} under ThreadSanitizer, is the program race-free? What can actually go wrong? What does mode \texttt{B} change? Four threads each hold their own copy of one \texttt{shared\_ptr} and copy it in a loop: is that a race? On g++ 11.4, can you use \texttt{std::atomic<std::shared\_ptr<int>{}>}?
\omcode{code/interviews-3/03-cpp/cpp/shared_race.cpp}{1}{25}{Replacing a shared pointer under a reader.}\label{lst:iv3:cp:sharedrace}
\end{question}

\begin{question}[$\star\star\star$ \iqroles{QD, Dev, LLE} \iqfirm{any} \iqlevel{junior, experienced}]\label{iq:iv3:cpp:51}
In C++20, write a bounded blocking queue for several producers and consumers, with \texttt{push}, \texttt{pop} and \texttt{close}. \texttt{push} blocks while full and returns \texttt{false} once closed. \texttt{pop} blocks while empty and returns an empty \texttt{std::optional} once closed and drained. Why two condition variables? Why does each wait take a predicate? What must \texttt{close} do so that nobody sleeps forever? How would you test it?
\end{question}

\begin{question}[$\star\star\star$ \iqroles{QD, Dev, LLE, QR} \iqfirm{any} \iqlevel{junior, experienced}]\label{iq:iv3:cpp:52}
In C++20, write \texttt{choose(n, k)} for the binomial coefficient, on \texttt{std::uint64\_t}, returning an empty \texttt{std::optional} exactly when the true answer does not fit in 64 bits. It must not report overflow for an answer that fits, even if a naive intermediate product would overflow. Why does the factorial formula fail, and from which $n$? Prove that your intermediate values never overflow unless the answer does. What is the smallest $n$ at which the answer can overflow?
\end{question}

\begin{question}[$\star\star\star$ \iqroles{QD, Dev, LLE} \iqfirm{any} \iqlevel{junior, experienced}]\label{iq:iv3:cpp:53}
\Cref{lst:iv3:cp:terminate} is run three times with g++ 11.4, with the argument \texttt{noexcept}, then \texttt{destructor}, then \texttt{plain}. What happens each time? Why does the \texttt{try} block not catch the first two? \texttt{Closer}'s destructor has no \texttt{noexcept} written on it, so why does it behave like the first? Could you let a destructor throw, and when would it still end the program?
\omcode{code/interviews-3/03-cpp/cpp/terminate.cpp}{2}{20}{Three ways to throw.}\label{lst:iv3:cp:terminate}
\end{question}

\begin{question}[$\star\star\star$ \iqroles{QD, Dev, LLE} \iqfirm{crypto, market maker, hedge fund} \iqlevel{junior, experienced}]\label{iq:iv3:cpp:54}
What does \cref{lst:iv3:cp:crtp} print? Explain how \texttt{Strat<Momo>} reaches \texttt{Momo::signal} with no virtual function. With g++ 11.4, \texttt{-O2}, what does the assembly of \texttt{static\_call} contain that differs from \texttt{virtual\_call}? Where does the speed really come from? What happens if \texttt{Momo} has no \texttt{signal}? Why can you not put a \texttt{Momo} and another strategy of this form into one container?
\omcode{code/interviews-3/03-cpp/cpp/crtp.cpp}{1}{31}{Static and dynamic dispatch.}\label{lst:iv3:cp:crtp}
\end{question}

\begin{question}[$\star\star\star$ \iqroles{QD, Dev, LLE} \iqfirm{any} \iqlevel{junior, experienced}]\label{iq:iv3:cpp:55}
Run with no arguments, what does \cref{lst:iv3:cp:constexpr} print? Replace the comment with each of \texttt{int d = twice(argc);}, \texttt{int e = always(argc);}, \texttt{constexpr int f = twice(argc);} and \texttt{constexpr int o = twice(2000000000);}. Which compile? How do \texttt{constexpr} and \texttt{consteval} on a function differ? The last line would be undefined at run time, so why is it a compile error here?
\omcode{code/interviews-3/03-cpp/cpp/constexpr_table.cpp}{1}{27}{Work at compile time.}\label{lst:iv3:cp:constexpr}
\end{question}

\begin{question}[$\star\star\star\star$ \iqroles{QD, Dev, LLE} \iqfirm{any} \iqlevel{experienced}]\label{iq:iv3:cpp:56}
In C++20, implement a \texttt{Vec<T>} with the destructor, copy and move, \texttt{emplace\_back}, \texttt{pop\_back} and \texttt{operator[]}, and which works for a type with no default constructor. \texttt{emplace\_back} must give the strong guarantee, and \texttt{v.emplace\_back(v[0])} must work when it triggers a reallocation. Why not \texttt{new T[c]}? In what order do the steps of a reallocation go? When do you move the old elements, and when copy them? How would you test the guarantee?
\end{question}

\begin{question}[$\star\star\star\star$ \iqroles{QD, Dev, LLE} \iqfirm{crypto, market maker, hedge fund} \iqlevel{experienced}]\label{iq:iv3:cpp:57}
In C++20, implement a reference-counted pointer \texttt{Rc<T>}, a \texttt{shared\_ptr} without weak pointers, with one allocation for count and object, and a count that is safe when several threads hold copies. Which memory order does the increment need? Which does the decrement need, and what breaks if it is relaxed? Show a test that would catch that on a strongly ordered x86-64. What does your class still not make safe?
\end{question}

\begin{question}[$\star\star\star\star$ \iqroles{QD, Dev, LLE} \iqfirm{crypto, market maker, hedge fund} \iqlevel{experienced}]\label{iq:iv3:cpp:58}
In C++20, implement \texttt{Fn<R(A.{}.{}.)>}: a move-only, type-erased callable with a 32-byte small buffer, falling back to the heap for a larger or a throwing-move callable. It must hold a lambda that captures a \texttt{std::unique\_ptr}. How do you call a function whose type you have forgotten? What do you need beyond the call? Why must the buffer be used only for a callable with a \texttt{noexcept} move? Which allocations are left? What does \texttt{std::function} have that your class lacks?
\end{question}

\begin{question}[$\star\star\star\star$ \iqroles{QD, Dev, LLE} \iqfirm{any} \iqlevel{experienced}]\label{iq:iv3:cpp:59}
In C++20, implement \texttt{Var<Ts.{}.{}.>}, a simplified \texttt{std::variant}, with \texttt{emplace<T>}, \texttt{get\_if<T>}, \texttt{index}, \texttt{visit} and the destructor; copying may be deleted. How big must the storage be and how must it be aligned, for any list of types? How does the destructor know what to destroy? What state are you in if a constructor throws inside \texttt{emplace}? How do you find the index of a type at compile time, with no recursion?
\end{question}

\begin{question}[$\star\star\star\star$ \iqroles{QD, Dev, LLE, QR} \iqfirm{any} \iqlevel{experienced}]\label{iq:iv3:cpp:60}
\texttt{d = a + b + c} on vectors of 1000 doubles, with an ordinary \texttt{operator+} returning a \texttt{std::vector<double>}, allocates twice and runs two addition loops. In C++20, write expression templates so that it allocates nothing, given that \texttt{d} already has its size, and runs one loop. What does \texttt{a + b} return now? Is \texttt{a = a + b} safe? Now \texttt{auto e = a + b + c; use(e[0]);}: what happens, and why did \texttt{auto} cause it? How would you repair the class?
\end{question}
